\section{Handling of Binders, Names, and Contexts}
\label{sec:binders-names-contexts}

Encoding a quantified or functional B construct forces the encoder to invent SMT
names for the variables it binds and to record their types. Both tasks are realized
as effects of the {\small\leancodeptr{Encoder}{Encoder/Basic.html\#Encoder}} monad
(\Cref{sec:lean-implementation-architecture}), acting on two components of its state
that are given deliberately different lifetimes.

\paragraph{Fresh names and the typing context.}
The operation {\small\leancodeptr{freshVar}{Encoder/Basic.html\#SMT.freshVar}}, given a type
$\tau$, returns a name distinct from every name used so far, records it among the
used names, and inserts the binding into the running context~$\Gamma$;
{\small\leancodeptr{addToContext}{Encoder/Basic.html\#SMT.addToContext}} and
{\small\leancodeptr{eraseFromContext}{Encoder/Basic.html\#SMT.eraseFromContext}} insert and
remove bindings explicitly. The fresh-variable counter and the set of used names are
\emph{monotone}---never rolled back---which guarantees that every generated name is
globally unique across the whole run. The typing context~$\Gamma$, by contrast, is
\emph{scoped}: a bound variable is meaningful only while its body is being encoded,
so~$\Gamma$ is extended on entering a binder and restored on leaving it.

\paragraph{Erasing binder variables.}
\inlremark{Decide if we keep this.}
The interplay is visible in
{\small\leancodeptr{defaultSpecM}{Encoder/Loosening/Rules.html\#defaultSpecM}}, which builds
the predicate expressing that a term takes the canonical default value of its
type. At a function type $\alpha \toS \beta$ the default is taken
pointwise, $\forallS x : \alpha.\ \mathsf{dflt}_{\beta}\bigl(t(x)\bigr)$, which in the
encoder reads:
\begin{leanlisting}
| .fun α β, t => do
    let x    ← freshVar α                     -- inserts x : α into Γ
    let body ← defaultSpecM β (.app t (.var x))
    eraseFromContext x                        -- x occurs only under the ∀ below
    pure (.forall [x] [α] body)
\end{leanlisting}
The fresh $x$ is drawn solely to serve as the universal binder. Since
\leaninl{freshVar} has inserted $x : \alpha$ into~$\Gamma$, but $x$ occurs only under
the $\forallS$ about to be built, it must be erased from~$\Gamma$ before the
predicate is returned: otherwise the result would carry a spurious free constant~$x$
and fail to typecheck in the ambient context, where $x$ is a variable bound by the
quantifier rather than a global symbol.

The same discipline recurs throughout the encoder. Every case of
{\small\leancodeptr{loosenAux}{Encoder/Loosening/Rules.html\#loosenAux_prf}} erases the
$\existsS$- and $\lamS$-binders it introduces once their specification is assembled,
and the universal case of {\small\leancodeptr{encodeTerm}{Encoder/Encoder.html\#encodeTerm}}
additionally snapshots~$\Gamma$ and the emitted declarations before encoding the
body: helper declarations produced while translating under the binder are then
re-scoped \emph{inside} the quantifier---as guarded universals---rather than left as
global constants that would escape their intended scope.

\begin{note}
  Keeping the fresh-name counter monotone while rolling back~$\Gamma$ is a matter of
  correctness, not convenience. Restoring the counter on leaving a binder could
  re-issue a name already attached to an emitted declaration, conflating two
  distinct SMT symbols; keeping it monotone preserves the global freshness invariant
  on which the correctness proofs rely, at the negligible cost of a few unused
  indices.
\end{note}
